% Run GCT_KRON_20261120; GPT-6 Astra Pro; effort not exposed.
\section{An endpoint-gradient normal form for the permanent residue tensor}

\subsection*{O-031. A proved reduction with an unresolved final square block}
This section concerns the actual tensor in the conditional
reduction QA033. It proves a simpler gradient formula and
one uniform invertibility statement. It does not assert
the arbitrary-size permanent endpoint.

Use the binary base \(M(t)=tJ+\operatorname{diag}(a_i)\),
with distinct \(a_i\), and the kernel basis \(K_{ab}\)
from (RJ.9). Put
\[
 \Delta_{ab}=\prod_{j\notin\{a,b\}}(a_j+t),\qquad
 P_{ab}=(1+t\partial_t)\Delta_{ab},\qquad
 Q_{ab}=\Delta_{ab}
             +2t\sum_{a<k<b}\Delta_{abk}\quad(a<b).
                                                               \tag{O-031.1}
\]
The empty interval contributes zero. Let
\(\mathcal T(t^j)=j!t^j\), as in QA033.
For linear forms \(p_i\), set
\[
 d_i=2a_i\sum_{j\ne i}\frac{p_i+p_j}{a_j-a_i},
                         \qquad e_i=d_i-2p_i.
                                                               \tag{O-031.2}
\]
Let \(B_p\in K\) have off-diagonal entries \(p_i+p_j\),
and hence diagonal entries \(d_i\). Let \(B_q\) be the
skew matrix with entries \(q_i-q_j\). Finally let \(C\)
be the skew matrix with upper-triangular entries \(1\).
All these matrices belong to \(K\).

Write \(H_M(B,D)=D^2\operatorname{per}_m(M)[B,D]\).
For \(a<b\), abbreviate \(A=a_a,B=a_b,D=B-A\).
The full inverse-factorial gradients are
\[
 \begin{split}
 \mathcal T^{-1}H_M(B_p,K_{ab})
   =\mathcal T^{-1}H_M(B_p,K_{ba})
   &=-\frac{A+B}{D}(p_a-p_b)P_{ab}
       -\frac{B e_a-A e_b}{D}\Delta_{ab},\\
 \mathcal T^{-1}H_M(B_q+\zeta C,K_{ab})
   &=-(q_a-q_b)P_{ab}-\zeta Q_{ab},\\
 \mathcal T^{-1}H_M(B_q+\zeta C,K_{ba})
   &=(q_a-q_b)P_{ab}+\zeta Q_{ab}.
 \end{split}                                                   \tag{O-031.3}
\]
Here \(\zeta\) is any further linear form. In particular,
all nonlocal dependence in the symmetric sector occurs
through the two endpoint forms \(e_a,e_b\), and
\(Q_{ab}/\Delta\) has residues only at the interval
vertices \(a,a+1,\ldots,b\).

\paragraph{Proof of the symmetric formula by actual scaling identities.}
Let \(P=\operatorname{diag}(p_i)\). Then
\[
 B_p=t^{-1}(PM+MP)
             +\operatorname{diag}(e_i-2a_ip_i/t).
\]
For a diagonal matrix \(P\), row and column scaling of
the permanent gives
\[
 D\operatorname{per}(M)[PM+MP]
                           =2\operatorname{tr}(P)\operatorname{per}(M).
\]
Differentiate this identity in a direction \(D_0\in K\).
Since \(D\operatorname{per}(M)[D_0]=0\),
\[
 H_M(PM+MP,D_0)
                       =-D\operatorname{per}(M)[PD_0+D_0P].
\]
The kernel vector \(D_0=K_{ab}\) depends rationally on
the parameters. Differentiating its identically zero
first-jet equation in \(a_i\) gives
\[
 H_M(E_{ii},K_{ab})
                     =-D\operatorname{per}(M)[\partial_{a_i}K_{ab}].
\]
Consequently
\[
 H_M(B_p,K_{ab})
 =-D\operatorname{per}(M)\left[
       \frac{PK_{ab}+K_{ab}P}{t}
        +\sum_i(e_i-2a_ip_i/t)\partial_{a_i}K_{ab}\right].
                                                               \tag{O-031.4}
\]
Only the two endpoint diagonal coordinates occur in
the derivatives of \(K_{ab}\).
They are \(A/D,-B/D\), with
\[
 \partial_A(A/D)=B/D^2,\quad
 \partial_B(A/D)=-A/D^2,
\]
and opposite derivatives for the other coordinate.

By (RJ.4), the relevant first cofactors are
\[
 \mathcal T(t\Delta_{ab}),\quad
 \mathcal T((B+t)\Delta_{ab}),\quad
 \mathcal T((A+t)\Delta_{ab}).
\]
Substituting them in (O-031.4), the terms containing
\(t^{-1}\mathcal T(\Delta_{ab})\) cancel.
The remaining expression is
\[
 -\frac{A+B}{D}(p_a-p_b)
                \frac{\mathcal T(t\Delta_{ab})}{t}
       -\frac{B e_a-A e_b}{D}\mathcal T(\Delta_{ab}).
\]
The identity
\[
 \frac{\mathcal T(tf)}{t}
                       =\mathcal T\bigl((1+t\partial_t)f\bigr)
\]
proves the first line of (O-031.3). Transposition gives
the same expression for \(K_{ba}\).
The argument was made in the rational-function field
where \(t\ne0\); the final expressions are polynomials,
so the result holds without that restriction.

\paragraph{The skew formulas and their interval term.}
For \(Q=\operatorname{diag}(q_i)\), diagonal conjugation
invariance of the permanent implies
\[
 D\operatorname{per}(M)[QM-MQ]=0.
\]
Since \(QM-MQ=tB_q\), differentiation along \(K_{ab}\)
gives
\[
 H_M(B_q,K_{ab})
     =-\frac1tD\operatorname{per}(M)[QK_{ab}-K_{ab}Q]
     =-(q_a-q_b)\frac{\mathcal T(t\Delta_{ab})}{t}.
\]
This proves the coboundary part of the skew formula.

For completeness, an arbitrary skew matrix \(b\) has
zero diagonal and belongs to \(K\). Directly collecting
the pair and triangle terms of its quadratic permanent
coefficient gives
\[
 \begin{split}
 \mathcal T^{-1}\frac12H_M(b,b)
 ={}&-\sum_{i<j}b_{ij}^{\,2}
                     (1+t\partial_t)\Delta_{ij}\\
    &+t\sum_{i<j<k}
               (b_{ij}+b_{jk}-b_{ik})^2\Delta_{ijk}.
 \end{split}                                                   \tag{O-031.5}
\]
Indeed, pairs contribute
\(-b_{ij}^2\Delta_{ij}\); a triangle contributes
\[
 -2t(b_{ij}b_{ik}-b_{ij}b_{jk}+b_{ik}b_{jk})\Delta_{ijk}.
\]
Writing this triangle term as a cycle square minus the
three edge squares yields (O-031.5), because
\(\partial_t\Delta_{ij}=\sum_{k\ne i,j}\Delta_{ijk}\).
Disjoint edges cancel by skew symmetry, equivalently
by (RJ.16).

At \(b_{ij}=q_i-q_j+\zeta\) for \(i<j\), every oriented
triangle cycle equals \(\zeta\). Differentiating
(O-031.5) in an edge \(a<b\), triangles with that edge
as the outer edge have negative cycle derivative;
the other triangles have positive derivative.
Since \(K_{ab}\) has half of the corresponding skew
variation, its skew gradient is
\[
 -(q_a-q_b+\zeta)P_{ab}
 +t\zeta\left(\sum_{k<a}\Delta_{kab}
       -\sum_{a<k<b}\Delta_{akb}
       +\sum_{k>b}\Delta_{abk}\right).
\]
Using the derivative formula for \(\Delta_{ab}\)
reduces this to the second line of (O-031.3).
Transposition gives the last line.

\paragraph{A uniform pure-skew isomorphism.}
Set \(\zeta=0\), let \(q_m=0\), and regard
\(q_1,\ldots,q_{m-1}\) as independent variables.
The quadratic differential obtained by varying the
skew edges by linear forms in these \(q\)-variables is
an isomorphism for every \(m\geq2\).

To prove this, apply the invertible binary operator
\((1+t\partial_t)^{-1}\) to (O-031.3), then pass to
residues after dividing by \(\Delta\).
Apart from the nonzero scalar \(1/(a_b-a_a)\), an
edge \(\{a,b\}\) tensored with a linear form \(\ell\)
maps to
\[
                    (e_a-e_b)\otimes(q_a-q_b)\ell .
                                                               \tag{O-031.6}
\]
The source and target dimensions are both
\(\binom m2(m-1)\).

For the dual injectivity argument, let \(T_i\) be
symmetric bilinear forms on the span of the \(q_i\),
modulo adding the same form to every \(T_i\), and set
\(T_m=0\). The edge conditions are
\[
               (T_a-T_b)(q_a-q_b,\xi)=0
                              \quad\hbox{for every }\xi.
\]
Edges \(\{i,m\}\) give \(T_i(q_i,\xi)=0\).
Edges \(\{i,j\}\) then give
\(T_i(q_j,\xi)=-T_j(q_i,\xi)\).
Taking \(\xi=q_j\) kills every diagonal coefficient.
For three distinct indices \(i,j,k\), symmetry and
three applications of the preceding relation yield
\[
 T_i(q_j,q_k)=-T_j(q_i,q_k)
             =T_k(q_i,q_j)=-T_i(q_j,q_k).
\]
Characteristic zero forces the coefficient to vanish.
Thus every \(T_i\) is zero, proving the isomorphism.

After removing the displayed nonzero edge weights
and the invertible binary operators, the selected
incidence map is precisely the integral map of O-020
with \(R=m-1\). Its integral cokernel is
\[
             (\mathbb Z/2\mathbb Z)^{\binom{m-1}{3}}.
\]
The unit eliminations and the remaining three-index
Smith blocks are proved there. This describes this
selected incidence map; it does not establish failure
of the full permanent jet projection in characteristic
two.

\paragraph{The tested one-charge candidate and its exact remaining scope.}
Use \(2m-1\) variables
\[
       p_1,\ldots,p_m,\quad q_1,\ldots,q_{m-1},
                         \qquad q_m=0,\quad\zeta=p_m,
\]
and set the free off-diagonal coefficients to
\[
             b_{ij}=p_i+p_j+q_i-q_j
                                +\operatorname{sgn}(j-i)p_m.
                                                               \tag{O-031.7}
\]
Complete the diagonal by the first-jet kernel equation.
Every skew triangle cycle is now the single variable
\(p_m\). This adds one fixed skew direction, multiplied
by one linear form, to the row/column-potential embedding.

The exact modular script
\path{agents/obstructions/check_one_charge_residue.py}
compares (O-031.3) with the original raw tensor and
then forms the complete square Jacobian.
For \(a_i=i\), it verifies all
\(90,336,900,1980\) gradient entries in sizes
\(m=3,4,5,6\), respectively, and gives full square
ranks \(30,84,180,330\).
The determinant residues modulo \(1000003\) are
\[
                       274064,\quad560456,\quad706498,\quad157229.
\]
Every rational denominator is nonzero modulo that
prime, so these are valid finite nonvanishing
certificates. They do not imply arbitrary-size
invertibility. The remaining mixed and symmetric
blocks of the one-charge map must still be controlled
uniformly; the pure-skew isomorphism alone does not
complete QA033.

\paragraph{Geometric and integral scope.}
All identities are obtained from actual permanent
cofactors, diagonal group actions, and explicit
polynomial coefficient maps. Their comodule
interpretation is SGA 3, I, Proposition 4.7.2,
pp.~34--35, with the grading in Proposition 4.7.3,
pp.~35--36. The calculations have integral models
after the required factorials and parameter
differences are inverted. The uniform isomorphism
argument also uses \(2\) invertible. The all-copy
consequence for the actual permanent coordinate ring
is conditional on the full square-map nonvanishing,
as in QA033 and F-016; no general irreducible
multiplicity claim is inferred from these partial
block computations.

